% 第 3 章 构建图书馆管理系统
\renewcommand{\codeInset}{18bp}
本章我们研究一个图书馆管理系统。与前面几章一样，我们继续开发 C++ 类。不过，本章要开发的是一个更接近真实世界的系统。本章的图书馆系统可以供真实的图书馆使用。图书馆由一组图书和一组读者构成。图书记录哪些读者借阅或预约了它们，读者则记录自己借阅和预约了哪些图书。

基本思想是：图书馆保存一组图书和一组读者。每本书都标记为已借出或未借出；如果已借出，就会存放借阅该书的读者的身份编号（identity number）。此外，一本书还可以被一个或多个读者预约，因此每本书还保存一份预约该书的读者编号列表。它必须是列表而不是集合，因为图书要按照读者预约的先后顺序依次借给这些读者。

每位读者保存两个集合，分别存放他们已借阅和已预约图书的编号。这两种情况下我们都使用集合（set）而不是列表（list），因为借阅或预约图书的先后顺序并不重要。

本章将涵盖以下主题：

\begin{itemize}
\item 使用由图书和读者构成的类，它们组成一个以整数编号为键的小型数据库。
\item 使用标准输入输出流：我们写出图书和读者的信息，并提示用户进行输入。
\item 使用文件处理与流。图书和读者通过标准 C++ 文件流写出和读入。
\item 最后，我们使用 C++ 标准库中泛型的集合、映射（map）与列表类。
\end{itemize}

\section{Book 类}

我们有三个类：Book、Customer 和 Library：

\begin{itemize}
\item \texttt{Book} 类记录一本书的信息。每本书都有作者和书名，以及一个唯一的编号。
\item \texttt{Customer} 类记录一位读者的信息。每位读者都有名字和地址，以及一个唯一的编号。
\item \texttt{Library} 类负责图书馆的各项操作，比如添加和删除图书与读者、借出、归还和预约图书，以及列出图书和读者。
\end{itemize}

\texttt{main} 函数只是创建一个 \texttt{Library} 类的对象。

此外，每本书都保存着当前是否被借出的信息；如果被借出，还会存放借阅该书的读者的编号。每本书还保存一份预约列表。同样，每位读者保存着当前借阅和预约图书的集合。

\texttt{Book} 类有两个构造函数。第一个是默认构造函数，用于从文件读取图书；第二个用于向图书馆添加新书，它以作者的名字和书名作为参数。

\medskip\noindent{\ttfamily\bfseries Book.h}

\begin{codebox}
class Book {
  public:
    Book(void);
    Book(const string& author, const string& title);
\end{codebox}

\texttt{author} 和 \texttt{title} 方法只是返回图书的作者和书名：

\begin{codebox}
    const string& author(void) const { return m_author; }
    const string& title(void) const { return m_title; }
\end{codebox}

图书馆的图书可以写入文件，也可以从文件读回：

\begin{codebox}
    void read(ifstream& inStream);
    void write(ofstream& outStream) const;
\end{codebox}

一本书可以被借出、预约或归还，预约也可以被取消。注意，图书被借出或预约时，我们需要提供读者的编号；而归还图书时则不需要，因为 \texttt{Book} 类自己记录着当前借阅该书的读者：

\begin{codebox}
    void borrowBook(int customerId);
    int reserveBook(int customerId);
    void unreserveBookation(int customerId);
    void returnBook();
\end{codebox}

图书被借出时，会存放借阅该书的读者的编号；\texttt{bookId} 返回图书自身的编号：

\begin{codebox}
    int bookId(void) const { return m_bookId; }
\end{codebox}

\texttt{borrowed} 方法在图书当前被借出时返回 \texttt{true}。此时，\texttt{customerId} 返回借阅该书的读者的编号：

\begin{codebox}
    bool borrowed(void) const { return m_borrowed; }
    int customerId(void) const { return m_customerId; }
\end{codebox}

一本书可以被一列读者预约，\texttt{reservationList} 返回这份预约列表：

\begin{codebox}
    list<int>& reservationList(void) { return m_reservationList; }
\end{codebox}

\texttt{MaxBookId} 字段是静态的，也就是说它为该类的所有对象所共有：

\begin{codebox}
    static int MaxBookId;
\end{codebox}

输出流运算符写出图书的信息：

\begin{codebox}
    friend ostream& operator<<(ostream& outStream,
                               const Book& book);
\end{codebox}

图书被借出时，\texttt{m\_borrowed} 字段为 \texttt{true}。图书的编号和可能的借阅者存放在 \texttt{m\_bookId} 和 \texttt{m\_customerId} 中：

\begin{codebox}
    private:
      bool m_borrowed = false;
      int m_bookId, m_customerId;
\end{codebox}

图书的作者名和书名存放在 \texttt{m\_author} 和 \texttt{m\_title} 中：

\begin{codebox}
      string m_author, m_title;
\end{codebox}

不止一位读者可以预约同一本书。此时，他们的编号存放在 \texttt{m\_reservationList} 中。它是列表而不是集合，因为预约是按先后顺序存放的。当一本书被归还时，按预约顺序排在下一位的读者就会借到这本书：

\begin{codebox}
      list<int> m_reservationList;
      };
\end{codebox}

本章我们使用 C++ 标准库中泛型的集合、映射与列表类。它们的声明存放在 Set、Map 和 List 头文件中。集合与列表类保存的集合和列表与上一章我们自己实现的类似。映射是这样一种结构：每个值都由一个唯一的键来标识，以便提供快速访问。

\medskip\noindent{\ttfamily\bfseries Book.cpp}

\begin{codebox}
#include <Set>
#include <Map>
#include <List>
#include <String>
#include <FStream>
using namespace std;

#include "Book.h"
#include "Customer.h"
#include "Library.h"
\end{codebox}

由于 \texttt{MaxBookId} 是静态的，我们用双冒号（\texttt{::}）记法来初始化它。每个静态字段都需要在类定义之外初始化：

\begin{codebox}
int Book::MaxBookId = 0;
\end{codebox}

默认构造函数什么也不做。它用于从文件读取时。尽管如此，我们仍然必须有一个默认构造函数，才能创建 \texttt{Book} 类的对象：

\begin{codebox}
Book::Book(void) {
  // Empty.
}
\end{codebox}

创建一本新书时，会为它分配一个唯一的编号。该编号存放在 \texttt{MaxBookId} 中，每创建一个新 \texttt{Book} 对象它就加一：

\begin{codebox}
Book::Book(const string& author, const string& title)
 :m_bookId(++MaxBookId),
  m_author(author),
  m_title(title) {
  // Empty.
}
\end{codebox}

\subsection{写入图书}

图书以类似的方式写入流。不过这次不是 read 而是 write。它们的工作方式相似：

\begin{codebox}
void Book::write(ofstream& outStream) const {
  outStream.write((char*) &m_bookId, sizeof m_bookId);
\end{codebox}

读取字符串时我们用 getline 而不是流运算符，因为流运算符只能读取一个单词，而 getline 能读取多个单词。不过写入流时，我们就可以用流运算符了，作者名和书名由一个还是多个单词都没有关系：

\begin{codebox}
  outStream << m_author << endl;
  outStream << m_title << endl;

  outStream.write((char*) &m_borrowed, sizeof m_borrowed);
  outStream.write((char*) &m_customerId, sizeof m_customerId);
\end{codebox}

与读取时一样，我们先写入预约列表中预约的个数，然后再逐个写入预约者的编号：

\begin{codebox}
  { int reservationListSize = m_reservationList.size();
    outStream.write((char*) &reservationListSize,
                  sizeof reservationListSize);

    for (int customerId : m_reservationList) {
      outStream.write((char*) &customerId, sizeof customerId);
    }
  }
}
\end{codebox}

\subsection{读取图书}

从文件读取任何类型的值（字符串除外）时，我们都使用 read 方法，它读取固定数目的字节。\texttt{sizeof} 运算符给出 \texttt{m\_bookId} 字段以字节为单位的大小。\texttt{sizeof} 运算符也可以用来求某个类型的大小。例如，\texttt{sizeof(int)} 给出 \texttt{int} 类型值以字节为单位的大小。类型必须用圆括号括起来：

\begin{codebox}
void Book::read(ifstream& inStream) {
  inStream.read((char*) &m_bookId, sizeof m_bookId);
\end{codebox}

从文件读取字符串值时，我们使用 C++ 标准函数 getline 来读取作者名和书名。如果名字由多个单词组成，使用输入流运算符是行不通的。如果作者名或书名由多个单词组成，只有第一个单词会被读入，其余的单词都不会被读取：

\begin{codebox}
  getline(inStream, m_author);
  getline(inStream, m_title);
\end{codebox}

注意，\texttt{m\_borrowed} 字段的值我们也用 read 方法读取，尽管它存放的是 \texttt{bool} 类型而不是 \texttt{int}：

\begin{codebox}
  inStream.read((char*) &m_borrowed, sizeof m_borrowed);
  inStream.read((char*) &m_customerId, sizeof m_customerId);
\end{codebox}

读取预约列表时，我们先读出列表中预约的个数，然后再读出各位预约者的编号：

\begin{codebox}
  { int reservationListSize;
    inStream.read((char*) &reservationListSize,
                sizeof reservationListSize);

    for (int count = 0; count < reservationListSize; ++count) {
      int customerId;
      inStream.read((char*) &customerId, sizeof customerId);
      m_reservationList.push_back(customerId);
    }
  }
}
\end{codebox}

\subsection{图书的借出与预约}

图书被借出时，\texttt{m\_borrowed} 变为 \texttt{true}，\texttt{m\_customerId} 被设置为借阅该书的读者的编号：

\begin{codebox}
void Book::borrowBook(int customerId) {
  m_borrowed = true;
  m_customerId = customerId;
}
\end{codebox}

图书被预约时的情况稍有不同。虽然一本书一次只能被一位读者借阅，却可以被多位读者预约。读者的编号被加入 \texttt{m\_reservationList}。列表的大小会被返回，让调用者知道自己在预约列表中的位置：

\begin{codebox}
int Book::reserveBook(int customerId) {
  m_reservationList.push_back(customerId);
  return m_reservationList.size();
}
\end{codebox}

图书被归还时，我们只是把 \texttt{m\_borrowed} 设为 \texttt{false}。不需要把 \texttt{m\_customerId} 设置成任何特定的值，只要图书未被借出，它就无关紧要：

\begin{codebox}
void Book::returnBook() {
  m_borrowed = false;
}
\end{codebox}

读者可以把自己从预约列表中移除。此时，我们对 \texttt{m\_reservationList} 调用 remove：

\begin{codebox}
void Book::unreserveBookation(int customerId) {
  m_reservationList.remove(customerId);
}
\end{codebox}

\subsection{显示图书}

输出流运算符写出图书的书名和作者。如果图书被借出，就写出读者的名字；如果预约列表非空，就写出各位预约读者的名字：

\begin{codebox}
ostream& operator<<(ostream& outStream, const Book& book) {
  outStream << "\"" << book.m_title << "\" by " << book.m_author;
\end{codebox}

访问静态字段（比如 Library 中的 \texttt{s\_customerMap}）时，我们使用双冒号（\texttt{::}）记法：

\begin{codebox}
  if (book.m_borrowed) {
    outStream << endl << "  Borrowed by: "
              << Library::s_customerMap[book.m_customerId].name()
              << ".";
  }

  if (!book.m_reservationList.empty()) {
    outStream << endl << "  Reserved by: ";
    bool first = true;

    for (int customerId : book.m_reservationList) {
      outStream << (first ? "" : ",")
                << Library::s_customerMap[customerId].name();
      first = false;
    }

    outStream << ".";
  }

  return outStream;
}
\end{codebox}

\section{Customer 类}

\texttt{Customer} 类记录一位读者的信息。它保存着该读者当前借阅和预约图书的集合。

\medskip\noindent{\ttfamily\bfseries Customer.h}

\begin{codebox}
class Customer {
  public:
    Customer(void);
    Customer(const string& name, const string& address);

    void read(ifstream& inStream);
    void write(ofstream& outStream) const;

    void borrowBook(int bookId);
    void reserveBook(int bookId);
    void returnBook(int bookId);
    void unreserveBook(int bookId);
\end{codebox}

\texttt{hasBorrowed} 方法在该读者当前至少借有一本书时返回 \texttt{true}。在下一节的 \texttt{Library} 类中，当前借有图书的读者是无法被删除的：

\begin{codebox}
    bool hasBorrowed(void) const { return !m_loanSet.empty(); }

    const string& name(void) const {return m_name;}
    const string& address(void) const {return m_address;}
    int id(void) const {return m_customerId;}
\end{codebox}

与前面用过的 \texttt{Book} 类一样，我们使用静态字段 \texttt{MaxCustomerId} 来统计读者的编号。我们还使用输出流运算符写出读者的信息：

\begin{codebox}
    static int MaxCustomerId;
    friend ostream& operator<<(ostream& outStream,
                               const Customer& customer);
\end{codebox}

每位读者都有名字、地址和一个唯一的编号。集合 \texttt{m\_loanSet} 和 \texttt{m\_reservationSet} 存放着该读者当前借阅和预约的图书的编号。注意，我们使用集合而不是列表，因为借阅和预约图书的先后顺序并不重要：

\begin{codebox}
    private:
      int m_customerId;
      string m_name, m_address;
      set<int> m_loanSet, m_reservationSet;
  };
\end{codebox}

\medskip\noindent{\ttfamily\bfseries Customer.cpp}

\begin{codebox}
#include <Set>
#include <Map>
#include <List>
#include <String>
#include <FStream>
using namespace std;

#include "Book.h"
#include "Customer.h"
#include "Library.h"
\end{codebox}

由于 \texttt{MaxCustomerId} 是静态字段，它需要在类之外定义：

\begin{codebox}
int Customer::MaxCustomerId;
\end{codebox}

默认构造函数只在从文件加载对象时使用，因此不需要对字段进行初始化：

\begin{codebox}
Customer::Customer(void) {
  // Empty.
}
\end{codebox}

第二个构造函数用于创建新的读者对象。我们用 \texttt{MaxCustomerId} 字段来初始化读者的编号，同时还初始化读者的名字和地址：

\begin{codebox}
Customer::Customer(const string& name, const string& address)
 :m_customerId(++MaxCustomerId),
  m_name(name),
  m_address(address) {
  // Empty.
}
\end{codebox}

\subsection{从文件读取读者}

read 方法从文件流中读取一位读者的信息：

\begin{codebox}
void Customer::read(ifstream& inStream) {
  inStream.read((char*) &m_customerId, sizeof m_customerId);
\end{codebox}

与 \texttt{Book} 类的 read 方法一样，我们必须使用 getline 函数而不是输入流运算符，因为输入流运算符只会读入一个单词：

\begin{codebox}
  getline(inStream, m_name);
  getline(inStream, m_address);

  { int loanSetSize;
    inStream.read((char*) &loanSetSize, sizeof loanSetSize);

    for (int count = 0; count < loanSetSize; ++count) {
      int bookId;
      inStream.read((char*) &bookId, sizeof bookId);
      m_loanSet.insert(bookId);
    }
  }

  { int reservationListSize;
    inStream.read((char*) &reservationListSize,
                sizeof reservationListSize);

    for (int count = 0; count < reservationListSize; ++count) {
      int bookId;
      inStream.read((char*) &bookId, sizeof bookId);
      m_loanSet.insert(bookId);
    }
  }
}
\end{codebox}

\subsection{将读者写入文件}

write 方法以与前面 \texttt{Book} 类相同的方式把读者的信息写入流：

\begin{codebox}
void Customer::write(ofstream& outStream) const {
  outStream.write((char*) &m_customerId, sizeof m_customerId);

  outStream << m_name << endl;
  outStream << m_address << endl;
\end{codebox}

写入集合时，我们先写入集合的大小，然后再逐个写入集合中的各个值：

\begin{codebox}
  { int loanSetSize = m_loanSet.size();
    outStream.write((char*) &loanSetSize, sizeof loanSetSize);

    for (int bookId : m_loanSet) {
      outStream.write((char*) &bookId, sizeof bookId);
    }
  }

  { int reservationListSize = m_reservationSet.size();
    outStream.write((char*) &reservationListSize,
                  sizeof reservationListSize);

    for (int bookId : m_reservationSet) {
      outStream.write((char*) &bookId, sizeof bookId);
    }
  }
}
\end{codebox}

\subsection{读者借阅与预约图书}

读者借阅一本书时，这本书被插入该读者的借阅集合：

\begin{codebox}
void Customer::borrowBook(int bookId) {
  m_loanSet.insert(bookId);
}
\end{codebox}

同样，读者预约一本书时，这本书被插入该读者的预约集合：

\begin{codebox}
void Customer::reserveBook(int bookId) {
  m_reservationSet.insert(bookId);
}
\end{codebox}

读者归还或取消预约一本书时，这本书就从借阅集合或预约集合中移除：

\begin{codebox}
void Customer::returnBook(int bookId) {
  m_loanSet.erase(bookId);
}

void Customer::unreserveBook(int bookId) {
  m_reservationSet.erase(bookId);
}
\end{codebox}

\subsection{显示读者}

输出流运算符写出读者的名字和地址。如果读者借阅或预约了图书，这些图书也会被写出：

\begin{codebox}
ostream& operator<<(ostream& outStream, const Customer& customer){
  outStream << customer.m_customerId << ". " << customer.m_name
            << ", " << customer.m_address << ".";

  if (!customer.m_loanSet.empty()) {
    outStream << endl << "  Borrowed books: ";

    bool first = true;
    for (int bookId : customer.m_loanSet) {
      outStream << (first ? "" : ",")
                << Library::s_bookMap[bookId].author();
      first = false;
    }
 }

  if (!customer.m_reservationSet.empty()) {
    outStream << endl << "  Reserved books: ";

    bool first = true;
    for (int bookId : customer.m_reservationSet) {
      outStream << (first ? "" : ",")
                << Library::s_bookMap[bookId].title();
      first = false;
    }
  }

  return outStream;
}
\end{codebox}

\section{Library 类}

最后，\texttt{Library} 类负责处理图书馆本身。它围绕借出和归还图书执行一系列任务。

\medskip\noindent{\ttfamily\bfseries Library.h}

\begin{codebox}
class Library {
  public:
    Library();
  private:
    static string s_binaryPath;
\end{codebox}

\texttt{lookupBook} 方法按作者和书名查找一本书。如果找到，它返回 \texttt{true}；此时该书的信息（一个 \texttt{Book} 类的对象）会被复制到 \texttt{bookPtr} 所指的对象中：

\begin{codebox}
    bool lookupBook(const string& author, const string& title,
                    Book* bookPtr = nullptr);
\end{codebox}

同样，\texttt{lookupCustomer} 按名字和地址查找一位读者。如果找到该读者，就返回 \texttt{true}，并把信息复制到 \texttt{customerPtr} 所指的对象中：

\begin{codebox}
    bool lookupCustomer(const string& name, const string& address,
                        Customer* customerPtr = nullptr);
\end{codebox}

本章的应用程序围绕下面这些方法展开。它们执行图书馆系统的各项任务。每个方法都会提示用户进行输入，然后执行一项任务，比如借出或归还一本书。

下面的方法各自执行一项任务，即查找图书或读者的信息、添加或删除图书、列出图书、在图书馆中添加和删除图书，以及借出、预约和归还图书：

\begin{codebox}
    void addBook(void);
    void deleteBook(void);
    void listBooks(void);
    void addCustomer(void);
    void deleteCustomer(void);
    void listCustomers(void);
    void borrowBook(void);
    void reserveBook(void);
    void returnBook(void);
\end{codebox}

load 和 save 方法分别在执行的开始和结束时被调用：

\begin{codebox}
    void load();
    void save();
\end{codebox}

有两个映射分别保存图书馆的图书和读者。如前所述，映射是这样一种结构：每个值都由一个唯一的键来标识，以便提供快速访问。图书和读者的唯一编号就是键：

\begin{codebox}
  public:
    static map<int,Book> s_bookMap;
    static map<int,Customer> s_customerMap;
};
\end{codebox}

\medskip\noindent{\ttfamily\bfseries Library.cpp}

\begin{codebox}
#include <Set>
#include <Map>
#include <List>
#include <String>
#include <FStream>
#include <IOStream>
#include <Algorithm>
using namespace std;

#include "Book.h"
#include "Customer.h"
#include "Library.h"

map<int,Book> Library::s_bookMap;
map<int,Customer> Library::s_customerMap;
\end{codebox}

两次执行之间，图书馆的信息存放在硬盘上的 Library.bin 文件中。注意，字符串中我们用两个反斜线表示一个反斜线：第一个反斜线表明后面的字符是特殊字符，第二个反斜线则表明它就是一个反斜线：

\begin{codebox}
string Library::s_binaryPath("Library.bin");
\end{codebox}

构造函数加载图书馆的信息，显示一个菜单，然后不断循环，直到用户退出。在执行结束之前，图书馆的信息会被保存：

\begin{codebox}
Library::Library(void) {
\end{codebox}

在显示菜单之前，先从文件加载图书馆的信息（图书、读者、借阅和预约）：

\begin{codebox}
  load();
\end{codebox}

只要 quit 不为 \texttt{true}，while 语句就会一直循环下去。直到用户从菜单中选择 Quit 选项，它才会变成 \texttt{true}：

\begin{codebox}
  bool quit = false;
  while (!quit) {
    cout << "1. Add Book" << endl
         << "2. Delete Book" << endl
         << "3. List Books" << endl
         << "4. Add Customer" << endl
         << "5. Delete Customer" << endl
         << "6. List Customers" << endl
         << "7. Borrow Book" << endl
         << "8. Reserve Book" << endl
         << "9. Return Book" << endl
         << "0. Quit" << endl
         << ": ";
\end{codebox}

用户从控制台输入流（\texttt{cin}）输入一个整数值，它存放在 \texttt{choice} 中：

\begin{codebox}
    int choice;
    cin >> choice;
\end{codebox}

我们用 switch 语句来执行所请求的任务：

\begin{codebox}
    switch (choice) {
      case 1:
        addBook();
        break;

      case 2:
        deleteBook();
        break;

      case 3:
        listBooks();
        break;

      case 4:
        addCustomer();
        break;

      case 5:
        deleteCustomer();
        break;

      case 6:
        listCustomers();
        break;

      case 7:
        borrowBook();
        break;

      case 8:
        reserveBook();
        break;

      case 9:
        returnBook();
        break;

      case 0:
        quit = true;
        break;
    }

    cout << endl;
  }
\end{codebox}

在程序结束之前，图书馆的信息会被保存：

\begin{codebox}
    save();
}
\end{codebox}

\subsection{查找图书与读者}

\texttt{lookupBook} 方法遍历图书映射。如果存在指定作者和书名的图书，它返回 \texttt{true}。如果该书存在，只要指针不为空，就把它的信息复制到 \texttt{bookPtr} 参数所指的对象中并返回 \texttt{true}。如果该书不存在，就返回 \texttt{false}，并且不向对象复制任何信息：

\begin{codebox}
bool Library::lookupBook(const string& author,
        const string& title, Book* bookPtr /* = nullptr*/) {
  for (const pair<int,Book>& entry : s_bookMap) {
    const Book& book = entry.second;
\end{codebox}

注意，\texttt{bookPtr} 可能是 \texttt{nullptr}。此时只返回 \texttt{true}，不向 \texttt{bookPtr} 所指的对象写入任何信息：

\begin{codebox}
    if ((book.author() == author) && (book.title() == title)) {
      if (bookPtr != nullptr) {
        *bookPtr = book;
      }

      return true;
    }
  }

  return false;
}
\end{codebox}

同样，\texttt{lookupCustomer} 遍历读者映射；如果存在指定名字的读者，它返回 \texttt{true}，并把读者信息复制到一个 \texttt{Customer} 对象中：

\begin{codebox}
bool Library::lookupCustomer(const string& name,
       const string& address, Customer* customerPtr /*=nullptr*/){
  for (const pair<int,Customer>& entry : s_customerMap) {
    const Customer& customer = entry.second;
\end{codebox}

同样，在这种情况下，\texttt{customerPtr} 也可能是 \texttt{nullptr}。此时只返回 \texttt{true}。添加新读者时，我们想知道是否已经存在同名同地址的读者：

\begin{codebox}
    if ((customer.name() == name) &&
        (customer.address() == address)) {
      if (customerPtr != nullptr) {
        *customerPtr = customer;
      }

      return true;
    }
  }

  return false;
}
\end{codebox}

\subsection{添加图书}

\texttt{addBook} 方法提示用户输入新书的作者名和书名：

\begin{codebox}
void Library::addBook(void) {
  string author;
  cout << "Author: ";
  cin >> author;

  string title;
  cout << "Title: ";
  cin >> title;
\end{codebox}

如果指定作者和书名的图书已经存在，就显示一条错误消息：

\begin{codebox}
  if (lookupBook(author, title)) {
    cout << endl << "The book \"" << title << "\" by "
         << author << " already exists." << endl;
    return;
  }
\end{codebox}

如果该书尚不存在，我们就创建一个新的 \texttt{Book} 对象，并把它加入图书映射：

\begin{codebox}
  Book book(author, title);
  s_bookMap[book.bookId()] = book;
  cout << endl << "Added: " << book << endl;
}
\end{codebox}

\subsection{删除图书}

\texttt{deleteBook} 方法提示用户输入图书的作者名和书名，如果该书存在就删除它：

\begin{codebox}
void Library::deleteBook() {
  string author;
  cout << "Author: ";
  cin >> author;

  string title;
  cout << "Title: ";
  cin >> title;
\end{codebox}

如果该书不存在，就显示一条错误消息：

\begin{codebox}
  Book book;
  if (!lookupBook(author, title, &book)) {
    cout << endl << "There is no book \"" << title << "\" by "
         << "author " << author << "." << endl;
    return;
  }
\end{codebox}

删除一本书时，我们遍历所有读者，并对每位读者归还并取消预约这本书。我们对每一本书都这样做，以防该书已被读者借出或预约。在下一章中，我们将使用指针，让我们能够以更高效的方式归还和取消预约图书。

注意，当我们遍历映射并获取每个 \texttt{Customer} 对象后，需要在修改其字段值之后把它放回映射中：

\begin{codebox}
  for (pair<int,Customer> entry : s_customerMap) {
    Customer& customer = entry.second;
    customer.returnBook(book.bookId());
    customer.unreserveBook(book.bookId());
    s_customerMap[customer.id()] = customer;
  }
\end{codebox}

最后，在确认图书存在、并且已经归还和取消预约之后，我们把它从图书映射中移除：

\begin{codebox}
  s_bookMap.erase(book.bookId());
  cout << endl << "Deleted." << endl;
}
\end{codebox}

\subsection{列出图书}

\texttt{listBooks} 方法相当简单。首先，我们检查图书映射是否为空。如果为空，就写出「No books.」；如果不为空，就遍历该映射，对每一本书都把它的信息写到控制台输出流（\texttt{cout}）：

\begin{codebox}
void Library::listBooks(void) {
  if (s_bookMap.empty()) {
    cout << "No books." << endl;
    return;
  }

  for (const pair<int,Book>& entry : s_bookMap) {
    const Book& book = entry.second;
    cout << book << endl;
  }
}
\end{codebox}

\subsection{添加读者}

\texttt{addCustomer} 方法提示用户输入新读者的名字和地址：

\begin{codebox}
void Library::addCustomer(void) {
  string name;
  cout << "Name: ";
  cin >> name;

  string address;
  cout << "Address: ";
  cin >> address;
\end{codebox}

如果同名同地址的读者已经存在，就显示一条错误消息：

\begin{codebox}
  if (lookupCustomer(name, address)) {
    cout << endl << "A customer with name " << name
         << " and address " << address << " already exists."
         << endl;
    return;
  }
\end{codebox}

最后，我们创建一个新的 \texttt{Customer} 对象，并把它加入读者映射：

\begin{codebox}
  Customer customer(name, address);
  s_customerMap[customer.id()] = customer;
  cout << endl << "Added." << endl;
}
\end{codebox}

\subsection{删除读者}

\texttt{deleteCustomer} 方法删除指定的读者（如果存在）：

\begin{codebox}
void Library::deleteCustomer(void) {
  string name;
  cout << "Name: ";
  cin >> name;

  string address;
  cout << "Address: ";
  cin >> address;

  Customer customer;
  if (!lookupCustomer(name, address, &customer)) {
    cout << endl << "There is no customer with name " << name
         << " and address " << address << "." << endl;
    return;
  }
\end{codebox}

如果该读者至少借有一本书，必须先归还，才能删除该读者：

\begin{codebox}
  if (customer.hasBorrowed()) {
    cout << "Customer " << name << " has borrowed at least "
         << "one book and cannot be deleted." << endl;
    return;
  }
\end{codebox}

不过，如果该读者预约了图书，我们在删除读者之前只需取消这些预约：

\begin{codebox}
  for (pair<int,Book> entry : s_bookMap) {
    Book& book = entry.second;
    book.unreserveBookation(customer.id());
    s_bookMap[book.bookId()] = book;
  }

  cout << endl << "Deleted." << endl;
  s_customerMap.erase(customer.id());
}
\end{codebox}

\subsection{列出读者}

\texttt{listCustomers} 方法的工作方式与 \texttt{listBooks} 类似。如果没有读者，就写出「No customers.」；如果有读者，就把他们写到控制台输出流（\texttt{cout}）：

\begin{codebox}
void Library::listCustomers(void) {
  if (s_customerMap.empty()) {
    cout << "No customers." << endl;
    return;
  }

  for (const pair<int,Customer>& entry : s_customerMap) {
    const Customer& customer = entry.second;
    cout << customer << endl;
  }
}
\end{codebox}

\subsection{借出图书}

\texttt{borrowBook} 方法提示用户输入图书的作者名和书名：

\begin{codebox}
void Library::borrowBook(void) {
  string author;
  cout << "Author: ";
  cin >> author;

  string title;
  cout << "Title: ";
  cin >> title;
\end{codebox}

如果不存在指定作者和书名的图书，就显示一条错误消息：

\begin{codebox}
  Book book;
  if (!lookupBook(author, title, &book)) {
    cout << endl << "There is no book \"" << title << "\" by "
         << "author " << author << "." << endl;
    return;
  }
\end{codebox}

此外，如果该书已经被借出，也显示一条错误消息：

\begin{codebox}
  if (book.borrowed()) {
    cout << endl << "The book \"" << title << "\" by " << author
         << " has already been borrowed." << endl;
    return;
  }
\end{codebox}

然后我们提示用户输入读者的名字和地址：

\begin{codebox}
  string name;
  cout << "Customer name: ";
  cin >> name;

  string address;
  cout << "Adddress: ";
  cin >> address;
\end{codebox}

如果不存在指定名字和地址的读者，就显示一条错误消息：

\begin{codebox}
  Customer customer;
  if (!lookupCustomer(name, address, &customer)) {
    cout << endl << "There is no customer with name " << name
         << " and address " << address << "." << endl;
    return;
  }
\end{codebox}

不过，如果图书存在且尚未被借出，而读者也存在，我们就把这本书加入该读者的借阅集合，并把该书标记为由该读者借出：

\begin{codebox}
  book.borrowBook(customer.id());
  customer.borrowBook(book.bookId());
\end{codebox}

注意，修改 \texttt{Book} 和 \texttt{Customer} 对象之后，我们必须把它们放回各自的映射中。在下一章中，我们将使用一种更直接的指针处理方式：

\begin{codebox}
  s_bookMap[book.bookId()] = book;
  s_customerMap[customer.id()] = customer;
  cout << endl << "Borrowed." << endl;
}
\end{codebox}

\subsection{预约图书}

\texttt{reserveBook} 方法的工作方式与 \texttt{borrowBook} 相同。它提示用户输入图书的作者名和书名：

\begin{codebox}
void Library::reserveBook(void) {
  string author;
  cout << "Author: ";
  cin >> author;

  string title;
  cout << "Title: ";
  cin >> title;
\end{codebox}

与 \texttt{borrowBook} 的情况类似，我们检查指定作者和书名的图书是否存在：

\begin{codebox}
  Book book;
  if (!lookupBook(author, title, &book)) {
    cout << endl << "There is no book \"" << title << "\" by "
         << "author " << author << "." << endl;
    return;
  }
\end{codebox}

不过，与 \texttt{borrowBook} 相比有一个不同之处：图书必须已经被借出，才能被预约。如果尚未借出，预约它毫无意义。在这种情况下，用户应该直接借阅这本书：

\begin{codebox}
  if (!book.borrowed()) {
    cout << endl << "The book with author " << author
         << " and title \"" << title << "\" has not been "
         << "borrowed. Please borrow the book instead." << endl;
    return;
  }
\end{codebox}

如果图书存在且已被借出，我们就提示用户输入读者的名字和地址：

\begin{codebox}
  string name;
  cout << "Customer name: ";
  cin >> name;

  string address;
  cout << "Address: ";
  cin >> address;
\end{codebox}

如果该读者不存在，就显示一条错误消息：

\begin{codebox}
  Customer customer;
  if (!lookupCustomer(name, address, &customer)) {
    cout << endl << "No customer with name " << name
         << " and address " << address << " exists." << endl;
    return;
  }
\end{codebox}

此外，如果这本书已经由该读者借出，我们就显示一条错误消息：

\begin{codebox}
  if (book.customerId() == customer.id()) {
    cout << endl << "The book has already been borrowed by "
         << name << "." << endl;
    return;
  }
\end{codebox}

如果图书存在且已被借出、但借出者不是该读者，我们就把该读者加入这本书的预约列表，并把这本书加入该读者的预约集合：

\begin{codebox}
  customer.reserveBook(book.bookId());
  int position = book.reserveBook(customer.id());
\end{codebox}

同样，在这种情况下，我们也必须把 \texttt{Book} 和 \texttt{Customer} 对象放回各自的映射中：

\begin{codebox}
  s_bookMap[book.bookId()] = book;
  s_customerMap[customer.id()] = customer;
\end{codebox}

最后，我们写出该读者在预约列表中的位置：

\begin{codebox}
  cout << endl << position << "nd reserve." << endl;
}
\end{codebox}

\subsection{归还图书}

\texttt{returnBook} 方法提示用户输入图书的作者名和书名：

\begin{codebox}
void Library::returnBook(void) {
  string author;
  cout << "Author: ";
  cin >> author;

  string title;
  cout << "Title: ";
  cin >> title;
\end{codebox}

如果该书不存在，就显示一条错误消息：

\begin{codebox}
  Book book;
  if (!lookupBook(author, title, &book)) {
    cout << endl << "No book \"" << title
         << "\" by " << author << " exists." << endl;
    return;
  }
\end{codebox}

如果该书尚未被借出，就显示一条错误消息：

\begin{codebox}
  if (!book.borrowed()) {
    cout << endl << "The book \"" << title
         << "\" by " << author
         << "\" has not been borrowed." << endl;
    return;
  }
\end{codebox}

与前面描述的方法不同，这种情况我们不需要询问读者。我们直接归还这本书，然后在预约列表中查找下一位读者：

\begin{codebox}
  book.returnBook();
  cout << endl << "Returned." << endl;

  Customer customer = s_customerMap[book.customerId()];
  customer.returnBook(book.bookId());
  s_customerMap[customer.id()] = customer;
\end{codebox}

如果这本书已被预约，我们就找出预约列表中的第一位读者，把他们从预约列表中移除，并让他们借阅这本书：

\begin{codebox}
  list<int>& reservationList = book.reservationList();

  if (!reservationList.empty()) {
    int newCustomerId = reservationList.front();
    reservationList.erase(reservationList.begin());
    book.borrowBook(newCustomerId);

    Customer newCustomer = s_customerMap[newCustomerId];
    newCustomer.borrowBook(book.bookId());

    s_customerMap[newCustomerId] = newCustomer;
    cout << endl << "Borrowed by " << newCustomer.name() << endl;
  }

  s_bookMap[book.bookId()] = book;
}
\end{codebox}

\subsection{将图书馆信息保存到文件}

保存图书馆的信息时，我们首先打开文件：

\begin{codebox}
void Library::save() {
  ofstream outStream(s_binaryPath);
\end{codebox}

如果文件被正确打开，我们先写出图书的数目，然后通过对 \texttt{Book} 对象调用 write 来写出每本书的信息：

\begin{codebox}
  if (outStream) {
    int numberOfBooks = s_bookMap.size();
    outStream.write((char*) &numberOfBooks, sizeof numberOfBooks);

    for (const pair<int,Book>& entry : s_bookMap) {
      const Book& book = entry.second;
      book.write(outStream);
    }
\end{codebox}

同样，我们写出读者的数目，然后通过调用 write 写出每位读者的信息：

\begin{codebox}
    int numberOfCustomers = s_customerMap.size();
    outStream.write((char*) &numberOfCustomers,
                  sizeof numberOfCustomers);

    for (const pair<int,Customer>& entry : s_customerMap) {
      const Customer& customer = entry.second;
      customer.write(outStream);
    }
  }
}
\end{codebox}

\subsection{从文件加载图书馆信息}

从文件加载图书馆信息时，我们对 read 的使用方式也与前面相同。我们首先打开文件：

\begin{codebox}
void Library::load() {
  ifstream inStream(s_binaryPath);
\end{codebox}

我们读出图书的数目，然后读出每本书的信息：

\begin{codebox}
  if (inStream) {
    int numberOfBooks;
    inStream.read((char*) &numberOfBooks, sizeof numberOfBooks);
\end{codebox}

对每本书，我们创建一个新的 \texttt{Book} 对象，通过调用 read 读出它的信息，并把它加入图书映射。我们还计算静态字段 \texttt{MaxBookId} 的新值，即把它设置为它自身与该书编号两者的最大值：

\begin{codebox}
    for (int count = 0; count < numberOfBooks; ++count) {
      Book book;
      book.read(inStream);
      s_bookMap[book.bookId()] = book;
      Book::MaxBookId = max(Book::MaxBookId, book.bookId());
    }
\end{codebox}

同样，我们读出读者的数目，然后通过调用 read 读出每位读者的信息：

\begin{codebox}
    int numberOfCustomers;
    inStream.read((char*) &numberOfCustomers,
                sizeof numberOfCustomers);
\end{codebox}

对每位读者，我们创建一个 \texttt{Customer} 对象，从文件中读出它的信息，把它加入读者映射，并计算静态字段 \texttt{MaxCustomerId} 的新值：

\begin{codebox}
    for (int count = 0; count < numberOfCustomers; ++count) {
      Customer customer;
      customer.read(inStream);
      s_customerMap[customer.id()] = customer;
      Customer::MaxCustomerId =
        max(Customer::MaxCustomerId, customer.id());
    }
  }
}
\end{codebox}

\section{main 函数}

最后，我们来编写执行这个图书馆的 \texttt{main} 函数。它非常简单；唯一要做的事情就是实例化一个 \texttt{Library} 类的对象。然后构造函数会显示主菜单：

\medskip\noindent{\ttfamily\bfseries Main.cpp}

\begin{codebox}
#include <Set>
#include <Map>
#include <List>
#include <String>
#include <FStream>
#include <IOStream>
using namespace std;

#include "Book.h"
#include "Customer.h"
#include "Library.h"

void main(void) {
  Library();
}
\end{codebox}

\section{小结}

本章我们构建了一个由 \texttt{Book}、\texttt{Customer} 和 \texttt{Library} 类组成的图书馆管理系统。

\texttt{Book} 类保存一本书的信息。每个 \texttt{Book} 对象都持有一个唯一的编号。它还记录着借阅者（如果该书已被借出）以及一份预约列表。同样，\texttt{Customer} 类保存一位读者的信息。与图书一样，每位读者持有一个唯一的编号。每个 \texttt{Customer} 对象还保存着已借阅和已预约图书的集合。最后，\texttt{Library} 类提供一系列服务，比如添加和删除图书与读者、借出、归还和预约图书，以及列出图书和读者。

本章中，每本书和每位读者都有一个唯一的编号。下一章我们将再次研究这个图书馆系统，不过我们会省去编号，改用指针来处理。
